\section*{अध्याय \ref{ch:b2:coordination-complexes} --- संक्रमण धातुएँ और उपसहसंयोजन संकुल}

\begin{solution}{exo:b2:coordination-complexes:1}
\ce{Ti^3+} $d^1$, \ce{V^3+} $d^2$, \ce{Cr^3+} $d^3$, \ce{Mn^2+} $d^5$, \ce{Fe^3+} $d^5$, \ce{Co^2+}
$d^7$, \ce{Ni^2+} $d^8$, \ce{Zn^2+} $d^{10}$।
\end{solution}

\begin{solution}{exo:b2:coordination-complexes:2}
जल 1, en 2, एथेनडाइओएट 2, ईडीटीए 6, थायोसायनेट 1; थायोसायनेट उभयदंतुर है
(S या N से)।
\end{solution}

\begin{solution}{exo:b2:coordination-complexes:3}\emergencystretch=3em
टेट्राऐम्मीनकॉपर(II) सल्फ़ेट; पोटैशियम हेक्सासायनिडोफ़ेरेट(III);
टेट्राऐम्मीनडाइक्लोरिडोकोबाल्ट(III) क्लोराइड; टेट्राकार्बोनिलनिकेल(0)।
\end{solution}

\begin{solution}{exo:b2:coordination-complexes:4}
रैखिक; चतुष्फलकीय; वर्ग समतलीय ($d^8$, तीसरी श्रेणी); अष्टफलकीय।
\end{solution}

\begin{solution}{exo:b2:coordination-complexes:5}
दो: cis (दोनों क्लोराइड सन्निकट) और trans (आमने-सामने)। चतुष्फलक में
केवल एक होता।
\end{solution}

\begin{solution}{exo:b2:coordination-complexes:6}
\ce{Cr(CO)6} 18; \ce{Fe(CO)5} 18; \ce{Ni(CO)4} 18; \ce{V(CO)6}: $5 + 12 = 17$;
\ce{Mn2(CO)10} प्रति मैंगनीज़ 18। \ce{V(CO)6} नियम तोड़ता है: यह \omterm{def:b2:diatomic-mos:magnetism}{अनुचुंबकीय} है और
आसानी से \ce{[V(CO)6]-} में अपचित हो जाता है, जिसके 18 हैं।
\end{solution}

\begin{solution}{exo:b2:coordination-complexes:7}
फ़ेरोसीन: Fe(II) $d^6$ + 12 = 18। कोबाल्टोसीन: Co(II) $d^7$ + 12 = 19। कोबाल्टोसीन
अपना अतिरिक्त इलेक्ट्रॉन आसानी से खो देता है, और 18 इलेक्ट्रॉनों वाला कोबाल्टोसीनियम आयन \ce{[Co(C5H5)2]+} देता है।
\end{solution}

\begin{solution}{exo:b2:coordination-complexes:8}
एक ईडीटीए जल के छह अणुओं को प्रतिस्थापित करता है: \ce{[Ni(H2O)6]^2+ + edta^4- -> [Ni(edta)]^2- +
6H2O}, दो कणों से सात। अभिक्रिया की एन्ट्रॉपी प्रबल रूप से धनात्मक है, और
$K$ बहुत बड़ा। छह अलग-अलग लिगैंड जल के छह अणुओं को कणों की संख्या में
कोई वृद्धि किए बिना प्रतिस्थापित करते हैं।
\end{solution}

\begin{solution}{exo:b2:coordination-complexes:9}
नाइट्रिटो-$\kappa N$ (\ce{Co-NO2}, नाइट्रो समावयवी, पीला) और नाइट्रिटो-$\kappa O$
(\ce{Co-O-N=O}, नाइट्रिटो समावयवी, लाल)।
\end{solution}

\begin{solution}{exo:b2:coordination-complexes:10}
तीन: trans समावयवी (अकिरैल: दर्पण तल) और cis
समावयवी के दोनों प्रतिबिंबरूप (दोनों कीलेट और दोनों क्लोराइड एक ओर या दूसरी ओर लिपटते हैं)।
\end{solution}

\begin{solution}{exo:b2:coordination-complexes:11}
समतलीय षट्भुज: B स्थितियों 1,2 (ऑर्थो), 1,3 (मेटा), 1,4 (पैरा) पर: तीन समावयवी।
त्रिकोणीय प्रिज़्म: B किसी त्रिभुज की भुजा के अनुदिश, किसी ऊर्ध्वाधर कोर के अनुदिश, या किसी
आयताकार फलक के आर-पार: तीन। अष्टफलक दो देता है। अनेक प्रयासों और कई यौगिकों के बावजूद
दो से अधिक न मिलना शेष दोनों ज्यामितियों को खारिज करता है --- यदि कोई
तीसरा समावयवी अनदेखा न रह गया हो।
\end{solution}

\begin{solution}{exo:b2:coordination-complexes:12}
\ce{[RhCl(PPh3)3]}: Rh(I) $d^8$ + 4 × 2 = 16। \ce{[IrH(CO)(PPh3)3]}: Ir(I) $d^8$ + 5 × 2 =
18। 16 इलेक्ट्रॉनों वाला रोडियम संकुल दो और ग्रहण कर सकता है (ऑक्सीकारी योग,
\cref{ch:b2:catalytic-cycles})।
\end{solution}

\begin{solution}{pb:b2:coordination-complexes:1}
\textbf{1.} \omterm{def:b2:coordination-complexes:coordination-sphere}{उपसहसंयोजन मंडल} के बाहर के क्लोराइड, जो विलयन में मुक्त हैं।
\textbf{2.} A: 0; B: 1; C और D: 2।
\textbf{3.} A: $[\ce{Co(NH3)6}]^{3+}$ + 3 \ce{Cl-}: 4 आयन; B: 3; C, D: 2। ये मेल खाते हैं।
\textbf{4.} +3।
\textbf{5.} $d^6$।
\textbf{6.} हर एक में छह: 6 \ce{NH3}; 5 \ce{NH3} + 1 Cl; 4 \ce{NH3} + 2 Cl।
\textbf{7.} \ce{[Co(NH3)6]Cl3}; \ce{[CoCl(NH3)5]Cl2}; \ce{[CoCl2(NH3)4]Cl} (दो बार)।
\textbf{8.} हेक्साऐम्मीनकोबाल्ट(III) क्लोराइड; पेंटाऐम्मीनक्लोरिडोकोबाल्ट(III) क्लोराइड।
\textbf{9.} टेट्राऐम्मीनडाइक्लोरिडोकोबाल्ट(III)।
\textbf{10.} ज्यामितीय (cis और trans) समावयवी।
\textbf{11.} \ce{[CoCl3(NH3)3]}, एक उदासीन संकुल: कोई आयन नहीं, कोई क्लोराइड तुरंत
अवक्षेपित नहीं।
\textbf{12.} दो: फ़ैक और मेर।
\textbf{13.} cis: दोनों Cl $90^\circ$ पर; trans: $180^\circ$ पर।
\textbf{14.} तीन (षट्भुज पर ऑर्थो, मेटा, पैरा स्थितियाँ)।
\textbf{15.} तीन।
\textbf{16.} यह अष्टफलक का समर्थन करता है, जो तीनों में अकेला दो का पूर्वानुमान करता है।
\textbf{17.} नहीं: तीसरे समावयवी का न मिलना नकारात्मक साक्ष्य है। तीसरा समावयवी
अष्टफलक को खारिज कर देता।
\textbf{18.} वह जो कीलेट बनाता है: द्विदंतुर एथेनडाइओएट केवल दो
cis स्थितियों को फैला सकता है।
\textbf{19.} कोबाल्ट के चारों ओर तीन en कीलेट, हर एक दो cis स्थितियों को फैलाता हुआ।
\textbf{20.} दोनों में से कोई नहीं: न दर्पण तल, न प्रतिलोमन केंद्र।
\textbf{21.} दो प्रतिबिंबरूप, $\Delta$ और $\Lambda$।
\textbf{22.} यह कि छह लिगैंड तीन विमाओं में, किसी अष्टफलक की तरह, व्यवस्थित हैं;
समतलीय व्यवस्था अकिरैल आयन देती।
\textbf{23.} cis: किरैल (दो प्रतिबिंबरूप)। trans: अकिरैल।
\textbf{24.} इसने दिखाया कि प्रकाशिक सक्रियता सामान्यतः आण्विक ज्यामिति का गुण है,
कार्बन का नहीं, और अष्टफलक की पुष्टि की।
\textbf{25.} अष्टफलक \ce{[CoCl2(NH3)4]+} के $\boldsymbol{2}$ समावयवियों का पूर्वानुमान करता है;
समतलीय षट्भुज और त्रिकोणीय प्रिज़्म $\boldsymbol{3}$ का।
\end{solution}
